\section{Grundlagen}
\label{sec:grundlagen}
Das folgende Kapitel legt die theoretischen und technischen Grundlagen dar, die für das Verständnis der in dieser Arbeit vorgestellten Lösung erforderlich sind. Zunächst wird die übergeordnete Plattform \gls{btp} eingeführt, auf der sämtliche Komponenten betrieben werden. Anschließend werden das \gls{cap}, die SAP-Fiori-Designrichtlinien, \gls{spa}, das \gls{odata}-Protokoll sowie das Konzept der Multi-Target Applications erläutert.

\subsection{SAP Business Technology Platform}
\label{subsec:btp}
Die SAP Business Technology Platform (\gls{btp}) ist die zentrale Cloud-Plattform der SAP SE, die Unternehmen eine integrierte Umgebung für die Entwicklung, Integration und Erweiterung betriebswirtschaftlicher Anwendungen bereitstellt \citep{metzger2021sap}. Sie vereint mehrere strategische Leistungsbereiche unter einem gemeinsamen Dach: Datenbank- und Datenmanagement, Analytik, Anwendungsentwicklung, Integration sowie Künstliche Intelligenz \citep{sap2024btp}.

Ein wesentliches Merkmal der \gls{btp} ist ihr Multi-Cloud-Ansatz. Die Plattform kann auf den Infrastrukturen der Hyperscaler Amazon Web Services, Microsoft Azure und Google Cloud Platform betrieben werden. Dadurch erhalten Unternehmen die Flexibilität, ihre Anwendungen in der jeweils bevorzugten Cloud-Umgebung bereitzustellen, ohne an einen einzelnen Infrastrukturanbieter gebunden zu sein.

Aus Sicht der Anwendungsentwicklung stellt die \gls{btp} verschiedene Laufzeitumgebungen zur Verfügung \citep{leymann2011cloud}. Die in dieser Arbeit verwendete Umgebung basiert auf \gls{cf}, einer quelloffenen \gls{paas}-Technologie. \gls{cf} ermöglicht es Entwicklern, Anwendungen in verschiedenen Programmiersprachen bereitzustellen, wobei die Plattform die zugrunde liegende Infrastruktur -- einschließlich Skalierung, Netzwerk und Betriebssystemverwaltung -- abstrahiert. Neben der Laufzeitumgebung bietet die \gls{btp} eine Vielzahl von verwalteten Diensten an, darunter Datenbanken wie SAP HANA Cloud, Authentifizierungsdienste (\gls{xsuaa}), Destination Services für die Anbindung externer Systeme sowie Dienste für die Prozessautomatisierung.

Die Administration der \gls{btp} erfolgt über das SAP BTP Cockpit, eine webbasierte Verwaltungsoberfläche. Hier werden Subaccounts, Umgebungen, Service-Instanzen und Berechtigungen konfiguriert. Subaccounts bilden dabei die organisatorische Einheit, in der Ressourcen und Berechtigungen gebündelt werden.

\subsection{SAP Cloud Application Programming Model}
\label{subsec:cap}
Das SAP Cloud Application Programming Model (\gls{cap}) ist ein quelloffenes Framework der SAP, das speziell für die Entwicklung von Unternehmensanwendungen auf der \gls{btp} konzipiert wurde. Es verfolgt einen deklarativen Ansatz, bei dem fachliche Domänenmodelle und Servicedefinitionen im Mittelpunkt stehen, während technische Aspekte wie Datenbankanbindung, Protokollhandhabung und Authentifizierung durch das Framework bereitgestellt werden \citep{sap2024cap}.

Das zentrale Beschreibungsmittel von \gls{cap} sind die Core Data Services (\gls{cds}) \citep{herger2020cap}. Mit \gls{cds} werden sowohl die Datenmodelle als auch die Serviceschnittstellen in einer kompakten, domänenspezifischen Sprache definiert. Aus einer \gls{cds}-Modellierung erzeugt das Framework automatisch die zugehörigen Datenbankartefakte, \gls{odata}-Serviceschnittstellen sowie die benötigte Laufzeitinfrastruktur. Dieser generative Ansatz reduziert den manuellen Implementierungsaufwand erheblich und stellt gleichzeitig die Einhaltung von Standards sicher.

\gls{cap} bietet darüber hinaus eine Reihe eingebauter Funktionalitäten, die für betriebswirtschaftliche Anwendungen typischerweise benötigt werden. Dazu zählen die automatische Bereitstellung von \gls{odata}-V4-Schnittstellen, die Integration mit dem \gls{xsuaa}-Dienst zur Authentifizierung und Autorisierung, die Unterstützung von Draft Handling für die Zwischenspeicherung von Nutzereingaben, sowie Mechanismen für die Mandantenfähigkeit. Als Laufzeitumgebung unterstützt \gls{cap} sowohl Node.js als auch Java. In der vorliegenden Arbeit kommt die Node.js-Variante zum Einsatz, die sich durch kurze Entwicklungszyklen und eine enge Integration mit dem \gls{cds}-Werkzeugkasten auszeichnet.

Ein weiterer Vorteil von \gls{cap} ist das integrierte Entwicklungswerkzeug \texttt{cds watch}, das während der lokalen Entwicklung Dateiänderungen erkennt und den Dienst automatisch neu startet. In Verbindung mit einer SQLite-Datenbank als lokalem Ersatz für SAP HANA Cloud ermöglicht dies einen schnellen Entwicklungs- und Testzyklus ohne Cloud-Abhängigkeiten.

\subsection{SAP Fiori und Fiori Elements}
\label{subsec:fiori}
SAP Fiori bezeichnet die Designrichtlinien und das Technologierahmenwerk der SAP für moderne Benutzeroberflächen betriebswirtschaftlicher Anwendungen. Die Fiori-Designprinzipien definieren fünf zentrale Leitlinien: Anwendungen sollen \textit{rollenbasiert} (auf die Aufgaben des Nutzers zugeschnitten), \textit{adaptiv} (geräteübergreifend nutzbar), \textit{kohärent} (konsistent im Erscheinungsbild), \textit{einfach} (auf das Wesentliche fokussiert) und \textit{ansprechend} (eine positive Nutzererfahrung bietend) sein \citep{sap2024fiori}.

Fiori Elements ist ein darauf aufbauendes UI-Framework, das Benutzeroberflächen auf Grundlage von \gls{odata}-Metadaten und UI-Annotationen generiert. Der wesentliche Vorteil dieses Ansatzes besteht darin, dass keine individuelle Frontend-Entwicklung erforderlich ist. Stattdessen werden die gewünschten UI-Eigenschaften -- etwa Tabellenspalten, Filterfelder, Formularabschnitte und Navigationsstrukturen -- deklarativ über Annotationen im \gls{cds}-Modell oder in separaten Annotationsdateien definiert. Das Framework erzeugt daraus zur Laufzeit die vollständige Benutzeroberfläche.

Fiori Elements stellt verschiedene Seitentypen bereit, die als sogenannte Floorplans bezeichnet werden. Die in der Praxis am häufigsten eingesetzten Floorplans sind der \textit{List Report}, der eine tabellarische Übersicht mit Such- und Filterfunktionen anzeigt, die \textit{Object Page}, die eine detaillierte Ansicht eines einzelnen Datenobjekts mit Kopfbereich und inhaltlichen Abschnitten darstellt, sowie die \textit{Worklist}, die eine aufgabenbezogene Liste ohne erweiterte Filtermöglichkeiten bereitstellt. Diese Floorplans können miteinander kombiniert werden, um vollständige Anwendungsszenarien abzubilden.

Technisch basiert Fiori Elements auf dem Framework \gls{ui5}, dem SAP-eigenen JavaScript-Framework für die Entwicklung von Webanwendungen \citep{bince2019fiori}. \gls{ui5} implementiert das Model-View-Controller-Muster und stellt eine umfangreiche Bibliothek vorgefertigter UI-Steuerelemente bereit \citep{sap2024fiorielem}. Fiori Elements nutzt diese Steuerelemente intern und abstrahiert deren Konfiguration durch den annotationsbasierten Generierungsansatz.

\subsection{SAP Build Process Automation}
\label{subsec:sbpa}
SAP Build Process Automation (\gls{spa}) ist ein cloudbasierter Dienst auf der \gls{btp}, der die Modellierung, Automatisierung und Überwachung von Geschäftsprozessen ermöglicht \citep{freund2019bpmn}. Als Nachfolger des früheren SAP Workflow Management verfolgt \gls{spa} einen Low-Code- bzw. No-Code-Ansatz, der es auch Fachanwendern ohne tiefgreifende Programmierkenntnisse erlaubt, Prozesse zu gestalten \citep{sap2024sbpa}.

Die zentralen Konzepte von \gls{spa} umfassen Prozesse, Formulare, \gls{api}-Auslöser, Aktionen, Bedingungen und Entscheidungstabellen. Ein \textit{Prozess} bildet den übergeordneten Ablauf ab und orchestriert die einzelnen Schritte. \textit{Formulare} dienen der Darstellung von Aufgaben und Eingabemasken für menschliche Bearbeiter. Über \textit{\gls{api}-Auslöser} können Prozesse programmatisch aus externen Anwendungen heraus gestartet werden, beispielsweise über \gls{rest}-Aufrufe. \textit{Aktionen} ermöglichen den Aufruf externer Dienste und \gls{api}s innerhalb eines Prozesses. \textit{Bedingungen} steuern den Kontrollfluss, und \textit{Entscheidungstabellen} erlauben die regelbasierte Auswertung von Geschäftslogik.

Für die Bearbeitung zugewiesener Aufgaben steht den Endnutzern die Anwendung \textit{My Inbox} zur Verfügung. Über My Inbox können Genehmigungsanfragen eingesehen, genehmigt oder abgelehnt werden. Die Anwendung folgt den SAP-Fiori-Designrichtlinien und ist sowohl über den SAP Build Lobby als auch über das SAP BTP Cockpit erreichbar.

Ein wesentlicher Aspekt von \gls{spa} im Kontext dieser Arbeit ist die Möglichkeit, Prozesse über \gls{api}-Auslöser zu starten und über Rückruf-Aktionen (\textit{Callbacks}) den Status an die auslösende Anwendung zurückzumelden. Dies ermöglicht eine ereignisgesteuerte Architektur, bei der die \gls{cap}-Anwendung einen Genehmigungsprozess auslöst und anschließend über das Ergebnis der Genehmigung informiert wird.

\subsection{OData-Protokoll}
\label{subsec:odata}
Das Open Data Protocol (\gls{odata}) ist ein offener Standard der Organisation OASIS (Organization for the Advancement of Structured Information Standards) für den Aufbau und die Nutzung von \gls{rest}ful \gls{api}s. In der aktuellen Version 4.0 definiert \gls{odata} ein standardisiertes Protokoll für \gls{crud}-Operationen auf Datenressourcen über das HTTP-Protokoll \citep{oasis2014odata}.

Ein \gls{odata}-Dienst beschreibt sein Datenmodell über ein maschinenlesbares Metadaten-Dokument, das unter dem Pfad \texttt{\$metadata} abrufbar ist. Dieses Dokument enthält die Definition aller Entitätstypen, deren Eigenschaften, Beziehungen untereinander (sogenannte Navigationseigenschaften) sowie die verfügbaren Operationen. Clients können dieses Metadaten-Dokument nutzen, um die Struktur des Dienstes zur Laufzeit zu ermitteln und generische Benutzeroberflächen zu erzeugen -- ein Prinzip, das von Fiori Elements intensiv genutzt wird.

Das \gls{odata}-Protokoll folgt dabei den Prinzipien von \gls{rest}, wie sie grundlegend von \citet{tilkov2015rest} beschrieben werden. Für die Abfrage von Daten stellt \gls{odata} standardisierte Abfrageoptionen bereit. Mit \texttt{\$filter} können Ergebnismengen nach bestimmten Kriterien eingeschränkt werden, \texttt{\$expand} ermöglicht das Nachladen verknüpfter Entitäten in einer einzelnen Anfrage, \texttt{\$select} beschränkt die zurückgegebenen Eigenschaften auf eine definierte Teilmenge und \texttt{\$orderby} legt die Sortierreihenfolge fest. Diese Abfrageoptionen werden als URL-Parameter übergeben und erlauben es dem Client, präzise zu steuern, welche Daten in welchem Umfang vom Server abgerufen werden.

Im Zusammenspiel mit dem \gls{cap} ergibt sich ein besonderer Vorteil: Das Framework generiert aus den \gls{cds}-Modellen automatisch vollständige \gls{odata}-V4-Dienste, sodass die Einhaltung des Standards ohne manuellen Implementierungsaufwand gewährleistet ist. Die \gls{odata}-Schnittstelle bildet damit die zentrale Kommunikationsschicht zwischen der serverseitigen Geschäftslogik und der clientseitigen Benutzeroberfläche.

\subsection{Multi-Target Applications}
\label{subsec:mta}
Das Konzept der Multi-Target Applications (\gls{mta}) adressiert die Herausforderung, dass Cloud-Anwendungen typischerweise aus mehreren, voneinander abhängigen Modulen bestehen, die als eine logische Einheit bereitgestellt werden müssen. Eine \gls{mta} fasst verschiedene Anwendungsmodule -- etwa serverseitige Dienste, Datenbankbereitstellungen und Benutzeroberflächen -- in einem gemeinsamen Archiv zusammen und beschreibt deren Abhängigkeiten und Konfigurationen in einem zentralen Deskriptor \citep{sap2024mta}.

Der \gls{mta}-Deskriptor ist eine YAML-Datei mit dem Namen \texttt{mta.yaml}, die im Wurzelverzeichnis des Projekts liegt. In dieser Datei werden die einzelnen Module, ihre Typen (z.\,B. \texttt{nodejs} für serverseitige Anwendungen, \texttt{hdb} für \gls{hdi}-Datenbankcontainer oder \texttt{approuter.nodejs} für den Application Router), ihre Abhängigkeiten untereinander sowie die benötigten Plattformdienste deklariert. Durch die zentrale Beschreibung aller Bestandteile wird sichergestellt, dass bei der Bereitstellung sämtliche Module in der korrekten Reihenfolge und mit der richtigen Konfiguration installiert werden.

Für die Erstellung des Bereitstellungsarchivs kommt das MTA Build Tool (\texttt{mbt}) zum Einsatz, das aus dem Quellcode und dem Deskriptor ein \texttt{.mtar}-Archiv erzeugt. Dieses Archiv wird anschließend mit dem \gls{cf}-CLI-Plugin \texttt{cf deploy} auf die \gls{btp} übertragen. Der Deployer auf der Plattform interpretiert den Deskriptor, erstellt die benötigten Service-Instanzen, stellt die Module bereit und konfiguriert die Verbindungen zwischen den Komponenten.

Im Kontext der vorliegenden Arbeit besteht die \gls{mta} aus drei Modulen: dem \gls{cap}-Serverdienst, der die Geschäftslogik und die \gls{odata}-Schnittstelle bereitstellt, dem \gls{hdi}-Datenbankbereitsteller, der die Datenbankartefakte in SAP HANA Cloud installiert, sowie dem Application Router, der die Fiori-Benutzeroberfläche ausliefert und die Authentifizierung über \gls{xsuaa} koordiniert. Ergänzt werden diese Module durch die konfigurierten Plattformdienste für Datenbank, Authentifizierung und Destination-Verwaltung.
